% =============================================================================
%              APPENDIX A: SOURCE CODE LISTINGS
% =============================================================================
\chapter{Source Code Listings}
\label{app:source_code}

This appendix provides the complete Python source code for all major components of the experimental pipeline described in Chapter~3. All experiments were implemented in Kaggle Notebooks using Python 3.12 and PyTorch.

\section{Data Acquisition}

\begin{lstlisting}[caption={FineGym clip download with \texttt{yt-dlp} temporal trimming.}, label={lst:fg_download}]
def download_clip(yt_id, start, end, out_path, max_retries=2):
    """Download a clip from YouTube using yt-dlp."""
    if os.path.exists(out_path) and \
       os.path.getsize(out_path) > 1000:
        return True # Already downloaded
    url = f'https://www.youtube.com/watch?v={yt_id}'
    for attempt in range(max_retries):
        try:
            cmd = [
                'yt-dlp', '--quiet', '--no-warnings',
                '--download-sections', f'*{start}-{end}',
                '-f', 'worst[ext=mp4]/worst',
                '--force-keyframes-at-cuts',
                '-o', out_path, url
            ]
            subprocess.run(
                cmd, capture_output=True, text=True, timeout=120)
            if os.path.exists(out_path) and \
               os.path.getsize(out_path) > 500:
                return True
        except (subprocess.TimeoutExpired, Exception):
            pass
    return False
\end{lstlisting}

\begin{lstlisting}[caption={HMDB51 download and local indexing pipeline.}, label={lst:hmdb_download}]
from huggingface_hub import hf_hub_download

# Download raw zip directly (bypasses schema errors)
archive_path = hf_hub_download(
    repo_id="jili5044/hmdb51",
    filename="hmdb51.zip",
    repo_type="dataset",
    local_dir=download_dir
)

# Extract and index videos by class folder
with zipfile.ZipFile(archive_path, 'r') as zip_ref:
    zip_ref.extractall(extract_dir)

# Build sample list from directory structure
ds = []
for root, dirs, files in os.walk(extract_dir):
    for file in files:
        if file.endswith(('.mp4', '.avi', '.mkv')):
            video_path = os.path.join(root, file)
            label_name = os.path.basename(root)
            ds.append({
                'video_path': video_path,
                'label': label_name
            })
\end{lstlisting}

\section{Feature Extraction}

\begin{lstlisting}[caption={V-JEPA 2 feature extraction with checkpointing.}, label={lst:vjepa_extraction}]
from transformers import VJEPA2Model, VJEPA2VideoProcessor

vproc = VJEPA2VideoProcessor.from_pretrained(
    'facebook/vjepa2-vith-fpc64-256')
vmodel = VJEPA2Model.from_pretrained(
    'facebook/vjepa2-vith-fpc64-256',
    torch_dtype=torch.float16,
    attn_implementation='sdpa',
).to(device).eval()

# Feature extraction with checkpointing
for i in range(start_idx, total):
    frames = load_frames(clip['path'], 8, 224)
    inputs = vproc(frames, return_tensors='pt')
    inputs = {k: v.to(device) for k, v in inputs.items()}
    with torch.no_grad():
        with torch.amp.autocast('cuda'):
            out = vmodel(**inputs)
            feat = (out.last_hidden_state.mean(dim=1)
                    .cpu().float().squeeze(0))
    features.append(feat)
\end{lstlisting}

\begin{lstlisting}[caption={Qwen2.5-VL 4-bit quantized feature extraction.}, label={lst:qwen_extraction}]
from transformers import (
    Qwen2_5_VLForConditionalGeneration,
    AutoProcessor, BitsAndBytesConfig)

bnb = BitsAndBytesConfig(
    load_in_4bit=True,
    bnb_4bit_compute_dtype=torch.float16,
    bnb_4bit_quant_type='nf4',
    bnb_4bit_use_double_quant=True)

qmodel = Qwen2_5_VLForConditionalGeneration.from_pretrained(
    'Qwen/Qwen2.5-VL-7B-Instruct',
    quantization_config=bnb,
    device_map='auto',
    torch_dtype=torch.float16)
qproc = AutoProcessor.from_pretrained(
    'Qwen/Qwen2.5-VL-7B-Instruct')

# Chat template construction for feature extraction
messages = [{'role': 'user', 'content': [
    {'type': 'video', 'video': frames},
    {'type': 'text', 'text': 'Describe.'},
]}]
text = qproc.apply_chat_template(
    messages, tokenize=False, add_generation_prompt=True)
inputs = qproc(text=[text], videos=[frames],
               return_tensors='pt', padding=True)

with torch.no_grad():
    out = qmodel(**inputs, output_hidden_states=True)
    feat = (out.hidden_states[-1].mean(dim=1)
            .cpu().float().squeeze(0))
\end{lstlisting}

\section{Classifier Architectures}

\begin{lstlisting}[caption={ResMLP architecture with residual blocks.}, label={lst:resmlp_arch}]
class ResBlock(nn.Module):
    def __init__(self, dim, dropout=0.2):
        super().__init__()
        self.net = nn.Sequential(
            nn.Linear(dim, dim),
            nn.BatchNorm1d(dim), nn.GELU(),
            nn.Dropout(dropout),
            nn.Linear(dim, dim),
            nn.BatchNorm1d(dim),
        )
        self.act = nn.GELU()

    def forward(self, x):
        return self.act(x + self.net(x))

class ResMLP(nn.Module):
    def __init__(self, in_dim, num_classes, hidden=256, n_blocks=3):
        super().__init__()
        self.proj = nn.Sequential(
            nn.Linear(in_dim, hidden),
            nn.BatchNorm1d(hidden), nn.GELU())
        self.blocks = nn.Sequential(
            *[ResBlock(hidden) for _ in range(n_blocks)])
        self.head = nn.Linear(hidden, num_classes)

    def forward(self, x):
        return self.head(self.blocks(self.proj(x)))
\end{lstlisting}

\begin{lstlisting}[caption={Training loop with label smoothing, mixup augmentation, and cosine LR scheduling.}, label={lst:training_loop}]
LABEL_SMOOTHING = 0.1

def mixup_data(x, y, alpha=0.2):
    """Feature-space mixup augmentation."""
    lam = np.random.beta(alpha, alpha)
    lam = max(lam, 1 - lam)
    idx = torch.randperm(x.size(0))
    return lam * x + (1-lam) * x[idx], y, y[idx], lam

def train_model(model, X_tr, y_tr, X_ev, y_ev,
                epochs=300, lr=1e-3, patience=30):
    opt = torch.optim.AdamW(
        model.parameters(), lr=lr, weight_decay=1e-4)
    scheduler = torch.optim.lr_scheduler.CosineAnnealingLR(
        opt, T_max=epochs)
    best_acc, best_state, wait = 0, None, 0

    for epoch in range(epochs):
        model.train()
        for xb, yb in dataloader:
            xb_mix, ya, yb_mix, lam = mixup_data(xb, yb)
            logits = model(xb_mix)
            loss_a = F.cross_entropy(
                logits, ya, label_smoothing=LABEL_SMOOTHING)
            loss_b = F.cross_entropy(
                logits, yb_mix, label_smoothing=LABEL_SMOOTHING)
            loss = lam * loss_a + (1 - lam) * loss_b
            opt.zero_grad()
            loss.backward()
            opt.step()
        scheduler.step()

        # Early stopping on eval accuracy
        model.eval()
        with torch.no_grad():
            preds = model(X_ev).argmax(1)
            acc = (preds == y_ev).float().mean().item()
        if acc > best_acc:
            best_acc = acc
            best_state = model.state_dict()
            wait = 0
        else:
            wait += 1
            if wait >= patience: break
    model.load_state_dict(best_state)
    return model
\end{lstlisting}

\section{HAF-STD Net Implementation}

\begin{lstlisting}[caption={HAF-STD Net architecture and spectral-temporal decomposition (HMDB51 configuration).}, label={lst:hafstd_arch}]
class HaarDWT1D(nn.Module):
    """1D Haar Discrete Wavelet Transform along temporal axis.
    Decomposes temporal signal into low-frequency (appearance/trend)
    and high-frequency (motion/change) components.
    """
    def forward(self, x):
        """x: (B, T, D) -> (low: (B, T//2, D), high: (B, T//2, D))"""
        if x.shape[1] % 2 != 0:
            x = F.pad(x, (0, 0, 0, 1))  # pad temporal dim if odd
        x_even = x[:, 0::2, :]  # Even-indexed temporal positions
        x_odd = x[:, 1::2, :]   # Odd-indexed temporal positions
        # Approximation = appearance/trend
        low = (x_even + x_odd) * 0.5
        # Detail = motion/change dynamics
        high = (x_even - x_odd) * 0.5
        return low, high


class HAFSTD_Net(nn.Module):
    """Hierarchical Adaptive Fusion with Spectral-Temporal Decomposition.
    Novel architecture for V-JEPA 2.1 on HMDB51 (51 classes).
    """
    def __init__(self, embed_dim=1280, num_layers=4, num_classes=51,
                 hidden=256, dropout=0.3):
        super().__init__()
        self.num_layers = num_layers
        # 11 streams: 4 appearance, 4 motion, 3 depth velocity
        self.num_streams = (num_layers * 2) + (num_layers - 1)  # 4*2 + 3 = 11
        self.dwt = HaarDWT1D()

        # Dedicated projection heads (11 total)
        self.stream_projections = nn.ModuleList([
            nn.Sequential(
                nn.Linear(embed_dim, hidden),
                nn.LayerNorm(hidden),
                nn.GELU(),
                nn.Dropout(dropout),
            )
            for _ in range(self.num_streams)
        ])

        # Adaptive routing network over 11 depth x frequency streams
        self.router = nn.Sequential(
            nn.Linear(self.num_streams * hidden, hidden),
            nn.GELU(),
            nn.Dropout(dropout),
            nn.Linear(hidden, self.num_streams),
        )

        # Classification head with cosine similarity
        self.final_proj = nn.Sequential(
            nn.Linear(hidden, hidden),
            nn.LayerNorm(hidden),
            nn.GELU(),
            nn.Dropout(dropout),
        )
        self.class_weights = nn.Parameter(
            torch.randn(num_classes, hidden))
        self.scale = nn.Parameter(torch.tensor(16.0))

    def forward(self, x):
        """x: (B, num_layers=4, T, D=1280)"""
        B = x.shape[0]
        streams = []

        # 1. Appearance & Motion streams (4 layers x 2 = 8 streams)
        for k in range(self.num_layers):
            layer_feat = x[:, k, :, :]  # (B, T, D)
            low, high = self.dwt(layer_feat)
            appearance = low.mean(dim=1)  # low-freq = appearance
            motion = high.mean(dim=1)     # high-freq = motion dynamics
            streams.append(self.stream_projections[k * 2](appearance))
            streams.append(self.stream_projections[k * 2 + 1](motion))

        # 2. Depth Representation Velocity streams (3 inter-layer deltas)
        for k in range(self.num_layers - 1):
            next_layer = x[:, k + 1, :, :].mean(dim=1)
            curr_layer = x[:, k, :, :].mean(dim=1)
            velocity = next_layer - curr_layer
            streams.append(self.stream_projections[8 + k](velocity))

        # Stack all 11 streams: (B, 11, hidden)
        all_streams = torch.stack(streams, dim=1)

        # 3. Adaptive Routing
        flat = all_streams.reshape(B, -1)
        routing_weights = torch.softmax(self.router(flat), dim=-1)
        fused = (all_streams * routing_weights.unsqueeze(-1)).sum(dim=1)

        # 4. Cosine Similarity Classification
        fused = self.final_proj(fused)
        fused_norm = F.normalize(fused, p=2, dim=-1)
        weight_norm = F.normalize(self.class_weights, p=2, dim=-1)
        logits = self.scale * fused_norm @ weight_norm.T
        return logits
\end{lstlisting}

\begin{lstlisting}[caption={4D Tensor Mixup and Training Harness for HAF-STD Net on HMDB51.}, label={lst:hafstd_train}]
def train_hafstd_model(model, X_tr_4d, y_tr, X_ev_4d, y_ev,
                       epochs=300, lr=1e-3, batch_size=64,
                       patience=30, use_mixup=True):
    """Train HAF-STD Net on 4D per-frame features (N, layers, T, D)."""
    model = model.to(device)
    opt = torch.optim.AdamW(
        model.parameters(), lr=lr, weight_decay=1e-4)
    scheduler = torch.optim.lr_scheduler.CosineAnnealingLR(
        opt, T_max=epochs)

    Xt = torch.tensor(X_tr_4d, dtype=torch.float32)
    yt = torch.tensor(y_tr, dtype=torch.long)
    Xe = torch.tensor(X_ev_4d, dtype=torch.float32).to(device)
    ye = torch.tensor(y_ev, dtype=torch.long).to(device)

    dl = DataLoader(
        TensorDataset(Xt, yt), batch_size=batch_size, shuffle=True)
    best_acc, best_state, wait = 0, None, 0

    for epoch in range(epochs):
        model.train()
        for xb, yb in dl:
            xb, yb = xb.to(device), yb.to(device)
            if use_mixup and xb.size(0) > 1:
                lam = max(np.random.beta(0.2, 0.2), 0.5)
                idx = torch.randperm(xb.size(0))
                xb_mix = lam * xb + (1 - lam) * xb[idx]
                logits = model(xb_mix)
                loss_a = F.cross_entropy(
                    logits, yb, label_smoothing=LABEL_SMOOTHING)
                loss_b = F.cross_entropy(
                    logits, yb[idx], label_smoothing=LABEL_SMOOTHING)
                loss = lam * loss_a + (1 - lam) * loss_b
            else:
                loss = F.cross_entropy(
                    model(xb), yb, label_smoothing=LABEL_SMOOTHING)
            opt.zero_grad(); loss.backward(); opt.step()
        scheduler.step()

        model.eval()
        with torch.no_grad():
            acc = (model(Xe).argmax(1) == ye).float().mean().item()
        if acc > best_acc:
            best_acc = acc
            best_state = {
                k: v.cpu().clone()
                for k, v in model.state_dict().items()
            }
            wait = 0
        else:
            wait += 1
            if wait >= patience: break
    model.load_state_dict(
        {k: v.to(device) for k, v in best_state.items()})
    return model
\end{lstlisting}
